\section{Conditional averaging in the shared environment}
\label{sec:averaging}

The table identities use conditional expectations on slices that retain
specified source and terminal vertices.  To approximate these expectations,
we need a walk that mixes on each slice.  We construct such a walk using
the sampling matrices of Section~\ref{sec:env}, then express its steps by
transformations that can be specified before the retained identifiers
are known.  Validation tests those identifiers when their vertices are
visited.

The identifiers are fixed when the conditional space is defined.  We
condition on one retained identifier per conditioned sampling matrix; additional
independent bit coordinates may be varied or frozen as specified below.

The \emph{full environment law} already includes the independent-column restriction
on each sampling matrix.  A conditional slice additionally imposes the indicated
retention events and fixes any designated frozen bits.  Its averaging
operator and $L_s$ norms use the normalized probability law on that slice.

The only external group-theoretic inputs are property~(T) of
$\mathrm{SL}_H(\mathbb Z)$ for $H\ge3$ and relative property~(T) of
\[
(\mathrm{SL}_2(\mathbb Z)\ltimes\mathbb Z^2,\mathbb Z^2).
\]
We use Shalom's quantitative statements:
\cite[Theorem~2.6, p.~156]{Shalom1999} for the linear group and
\cite[Theorem~2.1, p.~152]{Shalom1999} for the relative property.
We derive the affine finite-action gap and the conditional averaging walk from
these inputs below.

\begin{lemma}[A uniform affine spectral gap]\label{lem:env-affine}
Fix $H\ge3$ and write $G_H=\mathrm{SL}_H(\mathbb Z)\ltimes\mathbb Z^H$.
Let $\mathcal S$ consist of the elementary matrices with a single
$\pm1$ off the diagonal, acting linearly, and the translations by
$\pm e_i$, where $e_i$ are the standard basis vectors of $\mathbb Z^H$.
For the permutation representation $\rho$ of any finite transitive
action of $G_H$, with the uniform probability measure, the lazy
rational walk
\[
K_H=\frac12I+\frac1{2|\mathcal S|}\sum_{g\in\mathcal S}\rho(g)
\]
has a spectral gap bounded below by a positive constant $\gamma_H$
depending only on $H$.
\end{lemma}

\begin{proof}
The set $\mathcal S$ is a fixed finite symmetric generating set of $G_H$.
For a unitary representation $\rho$ and a vector $\xi$, put
\[
\delta=\max_{g\in\mathcal S}\|\rho(g)\xi-\xi\|.
\]
We first bound the distance from $\xi$ to the invariant subspace.

For each coordinate pair $\{i,j\}$, embed the two-dimensional affine
group on those coordinates, and let $P_{ij}$ be the orthogonal projection
onto the subspace fixed by its translations. This subspace and its
orthogonal complement are invariant under that pair group, because its
translation subgroup is normal. The eight generators in
\cite[Theorem~2.1, p.~152]{Shalom1999} belong to $\mathcal S$. Applying their
relative Kazhdan constant $1/10$ on the orthogonal complement gives
\[
 \frac1{10}\|\xi-P_{ij}\xi\|
 \le \max_{g\in\mathcal S_{ij}}
       \|\rho(g)(I-P_{ij})\xi-(I-P_{ij})\xi\|
 \le\delta,
\]
where $\mathcal S_{ij}$ is that eight-element subset. The last
inequality uses the commutation of $P_{ij}$ with the pair group
and the norm-one bound for $I-P_{ij}$. The first follows by
normalizing a nonzero vector in the complement: a smaller maximum
would give a translation-invariant vector in that complement,
contrary to its definition.
This is the projection argument in the proof of
\cite[Corollary~2.3, p.~154]{Shalom1999}.

To check commutation, let $U_i=\rho(e_i)$ and let $P_i$ project
onto $\ker(U_i-I)$. Each $U_j$ and its inverse preserve this
subspace because the translations commute; therefore $P_i$ commutes
with $U_j$. It then preserves both the fixed subspace of $U_j$ and
its orthogonal complement, so $P_iP_j=P_jP_i$. The product
$P_iP_j$ is the projection onto the vectors fixed by both coordinate
translations, hence equals $P_{ij}$. It follows that all $P_{ij}$
commute. Every coordinate occurs in some pair, so their product
$P_{\rm tr}$ projects onto the subspace fixed by all translations. With $q=\binom H2$, telescoping
the product gives
\[
\|\xi-P_{\rm tr}\xi\|
\le\sum_{i<j}\|\xi-P_{ij}\xi\|
\le10q\delta.
\]
Since the full translation subgroup is normal, $P_{\rm tr}$ commutes
with every $\rho(g)$. Its range carries a representation of
$G_H/\mathbb Z^H=\mathrm{SL}_H(\mathbb Z)$. Let $P_{\rm inv}$ project onto the
full $G_H$-invariant subspace, and let $\kappa_H>0$ be the Kazhdan
constant for the elementary generators from
\cite[Theorem~2.6, p.~156]{Shalom1999}; adjoining inverse generators
preserves the bound. Applying that constant on the range of
$P_{\rm tr}$ gives
\[
\|P_{\rm tr}\xi-P_{\rm inv}\xi\|\le\kappa_H^{-1}\delta,
\qquad
\|\xi-P_{\rm inv}\xi\|\le C_H\delta,
\quad C_H=10q+\kappa_H^{-1}.
\]
Thus the same generating set controls distance to invariants in every
unitary representation, with a constant depending only on $H$.

Now let $\rho$ be the permutation representation of a finite transitive
action, with the uniform probability measure, and put $m=|\mathcal S|$.
The operator $K_H$ of the statement is self-adjoint, and its spectrum lies in $[0,1]$. The invariant
functions are exactly the constants. For $f$ orthogonal to constants,
the preceding inequality and the Dirichlet identity give
\[
\langle f,(I-K_H)f\rangle
=\frac1{4m}\sum_{g\in\mathcal S}\|\rho(g)f-f\|_2^2
\ge\frac1{4mC_H^2}\|f\|_2^2.
\]
This proves the uniform spectral gap $\gamma_H=1/(4mC_H^2)$: every
eigenvalue of $K_H$ on functions orthogonal to constants is at most
$1-\gamma_H$.
\end{proof}

\paragraph{Coordinates and actions on a conditional slice.}
On a sampling matrix conditioned on retention of an identifier $y$ in a box of
side $L$, we use the coordinates $(F,w)$, where $F$ is the tuple of independent nonconstant columns of the
sampling matrix $A$ and $w=Av(y)\in\{0,\ldots,L-1\}^H$, identified with
an element of $(\mathbb Z/L\mathbb Z)^H$.  The constant column of $A$ is
uniquely determined by $F$, $w$ and $y$.  As observed in the proof of
Lemma~\ref{lem:env-ranks}, $F$ is uniform over independent column tuples and $w$ is
independent uniform in its box.  An element $(M,c)\in G_H$ acts by
\begin{equation}
 (F,w)\longmapsto(MF,\ Mw+c),
 \label{eq:env-slice-action}
\end{equation}
reducing modulo $P$ in the first coordinate and modulo $L$ in the
second.  On a sampling matrix without a conditioning identifier we use the same
description with a fixed public identifier and side $L=P$.
Concatenate the designated free bits into one block of length $H\nu$,
and identify it with $(\mathbb Z/2^\nu\mathbb Z)^H$, on which $(M,c)$ acts by
$v\mapsto Mv+c$ modulo $2^\nu$; frozen bits are not transformed.  The
proof of Lemma~\ref{lem:env-gap} shows that each of these actions is
transitive.

\begin{lemma}[Conditional averaging]\label{lem:env-gap}
Consider the fixed product of sampling matrices, conditioned on retention of
one specified identifier per conditioned sampling matrix. Include any designated
set of independent uniform bits, with length padded to a multiple of
$H$, and freeze any remaining independent bits.
Let $K$ be the operator on this conditional space that chooses one of
its boundedly many components uniformly, namely a sampling matrix or the block
of designated bits, applies one step of the walk $K_H$ of
Lemma~\ref{lem:env-affine} to that component through the action of
$G_H$ in Equation~\eqref{eq:env-slice-action}, or the affine action on
the free-bit block defined with it, and leaves the other
components unchanged. Then $K$ is a rational Markov averaging operator,
its stationary distribution is the indicated uniform conditional law
and, for each fixed $2\le s<\infty$,
\[
\|K^d-\Pi\|_{L_s\to L_s}\le C_s2^{-c_sd}
\qquad(d\ge0).
\]
Here $\Pi$ is conditional expectation over the variables being averaged.
The constants are independent of the field size, retention side
lengths, bit lengths, indicated identifiers, and frozen bits.
\end{lemma}

\begin{proof}
The sampling matrix action is transitive. Reduction of $\mathrm{SL}_H(\mathbb Z)$
contains the elementary generators of $\mathrm{SL}_H(\mathbb F_P)$.
Elementary row reduction generates $\mathrm{SL}_H(\mathbb F_P)$:
row additions are elementary matrices, while the needed determinant-one
rescalings are generated on two coordinates by
\[
 W(a)=
 \begin{pmatrix}1&a\\0&1\end{pmatrix}
 \begin{pmatrix}1&0\\-a^{-1}&1\end{pmatrix}
 \begin{pmatrix}1&a\\0&1\end{pmatrix}
 =\begin{pmatrix}0&a\\-a^{-1}&0\end{pmatrix},
 \qquad W(a)W(-1)=\diag(a,a^{-1}).
\]
The matrix $W(1)$ also supplies a row interchange with a sign.
Every elementary matrix over $\mathbb F_P$ lifts to an integer
elementary matrix. The finite group acts transitively on ordered tuples of independent
columns of length $r<H$: extend either column tuple to a basis and adjust
a complementary column to make the change-of-basis determinant one.
After choosing an integer lift $M$ that moves the column tuple, choose
$c$ with $c\equiv w'-Mw\pmod L$. Translations leave the column tuple
unchanged, so this reaches the desired pair $(F',w')$ even when
$P$ and $L$ differ. The action law holds because composition in
both coordinates is the reduction of
$(M,c)(M',c')=(MM',Mc'+c)$; no relation between the moduli is needed.
The affine action on the free-bit torus is transitive by translations.
There are a bounded number of components. Apply the lazy walk from
Lemma~\ref{lem:env-affine} to a uniformly selected component.
Let $K_i$ denote the single-component operator and let
$J_{\rm comp}$ be the number of components. The product operator is
$J_{\rm comp}^{-1}\sum_i K_i$, with $K_i$ acting as the identity
on other coordinates. A tensor product of eigenvectors has eigenvalue
$J_{\rm comp}^{-1}\sum_i\lambda_i$. All $\lambda_i$ lie in $[0,1]$;
unless every factor is constant, at least one is at most
$1-\gamma_H$. Thus the gap is at least
$\gamma_H/J_{\rm comp}$. A component consisting of one point
contributes only its constant eigenvector. The number of components
is bounded independently of the field and bit lengths, and every
$K_i$ preserves the product measure.

Writing its gap as $\gamma>0$, we have
$\|K^d-\Pi\|_{2\to2}\le(1-\gamma)^d$ and
$\|K^d-\Pi\|_{\infty\to\infty}\le2$, both with respect to the probability
law of the slice.  Interpolating between these two bounds with weight
$2/s$ on the $L_2$ bound, the convexity theorem of
Riesz~\cite[Theorems~II and~II$'$, p.~472]{Riesz1927} (the diagonal case of
the Riesz--Thorin theorem, which applies to this real operator) gives
\[
\|K^d-\Pi\|_{s\to s}
\le 2^{1-2/s}(1-\gamma)^{2d/s},
\]
which has the required form.
For Thorin's complex extension, see
\cite[Introduction, \S0.1; Russian translation, pp.~43--45]{Thorin1948}.
\end{proof}

\subsection{Generating paths without retaining their endpoints}

The sampling matrix coordinates for Lemma~\ref{lem:env-gap} use the indicated
identifier.  The next lemma is the implementation form of the same
averaging, and also supplies changes of variables for later error estimates.

\begin{lemma}[Validated walk variants]\label{lem:env-variants}
A length-$d$ word of the walk $K$ of Lemma~\ref{lem:env-gap}, that is, a
sequence of $d$ choices of a component and of a lazy step or a generator
in $\mathcal S$, can be expanded
into a bounded-alphabet family of variants described by $O(d)$ bits
such that:
\begin{enumerate}
\item each fixed variant acts bijectively and measure-preservingly on
the full environment before any endpoint tests;
\item its transformation is independent of the indicated identifiers;
\item for each walk word and each environment in the required
conditional slice, exactly one variant passes the slice and wrap
tests at the indicated identifiers, and this variant gives the
conditional walk transformation;
\item the sum of valid variants, each with the probability of its
original walk word, is exactly the conditional average.
\end{enumerate}
The tests can be performed when their indicated vertices are visited.
In the admissible implementation range, field elements and identifiers
have $O(B)$ bits, the total auxiliary bit string has $O(B)$ bits, and $d=O(B)$.
In that range the transformations and tests use $O(B)$ shared scratch space, and
the retained description of a variant uses $O(d)$ bits.
\end{lemma}

\begin{proof}
Consider a fixed affine generator $(M,c)$ on a sampling matrix with retention
side $L$. Enumerate the transformations
\begin{equation}\label{eq:env-wrap}
A'=MA+(c+Lz)e_1^{\mathsf T}\pmod P,
\end{equation}
where $e_1$ selects the constant column and $z$ ranges over a fixed
finite box of integer vectors. That box contains all possible wraps:
for $w\in\{0,\ldots,L-1\}^H$, every coordinate of $Mw+c$ is bounded
in absolute value by a fixed multiple of $L$ plus a fixed constant.
Since $L\ge1$ and the generators are fixed, one bound for $z$ works
for every side length.

Each transformation in Equation~\eqref{eq:env-wrap} is invertible.
Left multiplication preserves column independence, and the added
term changes only the constant column. It therefore preserves uniform
measure on the full sampling matrix. It does not refer to the indicated
identifier.

At that identifier, let $w=Av(y)$ in its integer representatives.
Check first that $w$ is in the retention box, and then that
$Mw+c+Lz$ lies in the same box coordinatewise, using integer arithmetic
before reduction modulo $P$. The unique valid wrap has coordinates
$z_i=-\lfloor(Mw+c)_i/L\rfloor$. For this vector the evaluation of
$A'$ is the representative of the intended affine action modulo $L$.
For a word with steps $(M_i,c_i,z_i)$, put
\[
 U_0=I,\quad v_0=0,\qquad
 U_i=M_iU_{i-1},\quad
 v_i=M_iv_{i-1}+c_i+Lz_i\pmod P.
\]
The full sampling matrix at step $i$ is
$A_i=U_iA+v_i e_1^{\mathsf T}$, so at every identifier
$A_iv(y)=U_iAv(y)+v_i$. The integer wrap test uses the
representative of $A_{i-1}v(y)$ and the current $c_i+Lz_i$
before reduction. Inductively it selects exactly the conditional
action modulo $L$ at each step. This proves uniqueness for the
whole word and also shows how to regenerate any intermediate
test from the word without saving its conditioning identifier.

The construction applies separately to the bounded number of sampling matrices.
On a sampling matrix of side $P$ the ordinary field reduction suffices;
alternatively, the same wrap validation may be used at its public
point. The free-bit torus has ordinary invertible affine operations
and requires no unknown endpoint. Frozen bits remain fixed.
Thus a word has only constantly many generator and wrap choices
per step, giving the $O(d)$-bit description.

For each word and initial environment in the slice, exactly one
variant survives. Its weight must be the weight of that word.
In particular, the sum is not divided by the number of guessed
wraps. This proves the asserted exact identity.
All matrices have fixed dimensions and $O(B)$-bit entries;
evaluating a bounded-degree polynomial or a word's individual steps
needs $O(B)$ arithmetic scratch. A test depending on a vertex can
be delayed until that vertex is current, with the required matrix
state regenerated from the recorded transformations.
Theorem~\ref{thm:ctrl-realization} accounts for reversible restoration
and the simultaneous space used when these operations are nested.
\end{proof}

\begin{remark}[Effective environment coordinates]
\label{rem:env-effective}
Conditional averaging is expressed in the sampling matrix and bit coordinates
of the environment supplied to that invocation. A fixed variant from
Lemma~\ref{lem:env-variants} is a measure-preserving change of variables
on the full environment; its retention and wrap tests remain explicit masks.
Thus free and frozen bits always refer to the supplied coordinates,
which may differ from those before the change of variables.
\end{remark}

\begin{remark}\label{rem:env-averaging}
Two distinct uses of the variants should be kept separate.
For centered-noise estimates, one first recognizes the sum of valid
variants as the conditional walk and applies
Lemma~\ref{lem:env-gap}. For crude bounds and changes of variables,
each fixed variant is a measure-preserving bijection, and one may
then pay for their $2^{O(d)}$ number. Counting variants before a
spectral estimate would needlessly weaken the useful contraction.
\end{remark}
